\begin{explanationbox}
	\textbf{\textcolor{CExplanation}{一句话直觉}}

	抽象函数定义域问题中，定义域始终是自变量$x$的取值范围，与复合形式无关。

	\medskip
	\textbf{\textcolor{CExplanation}{核心拆解}}
	\begin{description}[style=nextline, leftmargin=2.8em, labelsep=0.5em, itemsep=0.45em]
		\item[\textbf{要点一（定义域本质）：}] 定义域是自变量$x$的允许取值范围，不是中间变量的范围。
		\item[\textbf{要点二（变换规则）：}] 已知$f(x)$定义域求$f[g(x)]$定义域，解不等式；已知$f[g(x)]$定义域求$f(x)$定义域，求值域。
	\end{description}

	\medskip
	\textbf{\textcolor{CExplanation}{几何本质（投影关系）}}

	定义域对应函数图像在$x$轴上的投影区间。

	\medskip
	\textbf{\textcolor{CExplanation}{代数意义（变量代换）}}

	通过$t=g(x)$进行变量转换，将复合函数定义域问题转化为简单不等式或值域问题。

	\medskip
	\textbf{\textcolor{CExplanation}{考点价值（解题关键）}}
	\begin{itemize}[leftmargin=*, itemsep=0.5em, topsep=0.4em]
		\item \textbf{考法一（直接套用）：} 直接利用规则求复合函数定义域。
		\item \textbf{考法二（逆向推导）：} 借助规则反求参数或原函数定义域。
	\end{itemize}

	\medskip
	\textbf{\textcolor{CExplanation}{顿悟点}}

	\textbf{始终问自己：哪个是自变量$x$？不要被括号内表达式迷惑。}

	\medskip
	\textbf{\textcolor{CExplanation}{使用场景}}
	\begin{itemize}[leftmargin=*, itemsep=0.5em, topsep=0.4em]
		\item \textbf{场景一（抽象函数题）：} 如已知$f(x+1)$定义域求$f(2x)$定义域。
		\item \textbf{场景二（参数问题）：} 已知复合定义域限制，求参数范围。
	\end{itemize}
\end{explanationbox}
